\section{Uso de modelos preentrenados}\label{g-0}

Economicon no entrena modelos propios: usa modelos preentrenados a través de un gateway y los guía con instrucciones y contexto. La elección importa porque determina el coste por consulta, el tiempo de respuesta y dónde viajan los datos del cliente. Este apartado explica por qué se descartó entrenar, por qué ruta se accede a los modelos, cómo se eligió el modelo y qué se puede afirmar sobre su coste y su latencia. La tabla~\ref{tab-g1} resume los datos clave.

\begin{table}[!htbp]
\centering\small
\caption{Configuración de la versión final y evidencias de la elección}\label{tab-g1}
\begin{tabular}{@{}p{0.22\linewidth}p{0.3\linewidth}p{0.4\linewidth}@{}}
\toprule
\textbf{Ámbito} & \textbf{Dato o elección} & \textbf{Qué cuenta} \\
\midrule
Modelo de chat por defecto & DeepSeek V4.1 Flash; proveedor DeepInfra, 8 bits; sin ruta de reserva & Único generador de la configuración versionada; no certifica qué modelo respondió en cada llamada. \cite{config-gateway,arquitectura-final} \\
Modelo de embeddings & text-embedding-3-small & Genera los embeddings de la recuperación; no genera respuestas. \cite{config-gateway} \\
Alternativa nativa en AWS & Amazon Nova Micro y Titan V2 & Implementada y probada sin conexión; apagada por defecto y sin invocaciones reales. \cite{aws-ensayo} \\
Comparativa de contenido & 27 de 28 respuestas aceptables con cada modelo & GLM-5.2 frente a Flash con las mismas preguntas; valoración de un juez IA. \cite{comparativa-modelos} \\
Comparativa de coste & GLM-5.2: 0,1806 USD; Flash: 0,0392 USD & Las 28 operaciones de cada modelo, con los embeddings de las preguntas; sin la indexación inicial. \cite{comparativa-modelos} \\
Coste facturado por consulta & GLM-5.2: 0,0026 USD; Flash: 0,0006 USD & Consultas de chat terminadas, contrastadas con la factura del proveedor; ventanas de medición distintas. \cite{jup130-coste,jup133-coste} \\
Latencia histórica del banco & Percentil 95: 41,3 s antes del arreglo y 42,4 s tras él; objetivo de 10 s & Tiempo bajo el que quedan el 95\,\% de las operaciones completadas del banco de 28 preguntas. \cite{reevaluacion-modelo,banco28-tras} \\
\bottomrule
\end{tabular}
\end{table}

\subsection{Adaptación sin entrenamiento propio}\label{g-1}

Para un modelo preentrenado caben tres opciones: entrenarlo desde cero, ajustar sus pesos con datos propios (fine-tuning) o usarlo tal cual con instrucciones y contexto. Para el MVP se eligió la tercera, que concentra el esfuerzo en lo que diferencia a Economicon y evita mantener un ciclo propio de entrenamiento: las fuentes, los controles de la respuesta y el cálculo determinista de las cifras. El modelo de embeddings también se usa sin modificar. \cite{especificaciones,config-gateway}

\subsection{Ruta y condiciones de uso}\label{g-2}

LiteLLM actúa como gateway de modelos: una entrada común que guarda las credenciales, aplica límites de gasto por clave y registra cada consulta. OpenRouter encamina después la petición al proveedor. La aplicación solo conoce un alias estable, de modo que cambiar de modelo es un cambio de configuración y no de código. La decisión sobre el modelo por defecto sigue en estado de propuesta, a la espera de su ratificación por el equipo. \cite{config-gateway,adr-registro}

La configuración exige retención cero de datos (ZDR, Zero Data Retention), excluye a los proveedores que recopilan los datos de las peticiones y desactiva el cambio automático de proveedor ante un fallo. La cuantización a 8 bits representa los parámetros del modelo con menor precisión para abaratar la inferencia. Son condiciones configuradas, no una atestación externa de cada llamada. \cite{config-gateway}

Para el ensayo en AWS se implementó una ruta alternativa con los modelos nativos de la nube, sin pasar por el gateway. El código y sus controles están probados sin conexión, pero no se ha invocado ningún modelo de AWS, así que esta memoria no aporta mediciones de esa ruta. \cite{aws-ensayo}

\subsection{Cómo se eligió el modelo}\label{g-3}

GLM-5.2 fue el primer modelo del chat. Para sustituirlo se compararon GLM-5.2 y DeepSeek V4.1 Flash con las mismas preguntas, el mismo corpus y el mismo contexto. Ambos obtuvieron casi todas las respuestas aceptables según el juez IA; Flash costó menos de la cuarta parte; GLM-5.2 necesitó menos reintentos y tuvo tiempos más regulares. Pesó el coste, porque la calidad observada era equivalente. La valoración del juez describe aquella muestra y no sustituye una validación humana. \cite{comparativa-modelos}

Una evaluación posterior usó reglas automáticas y una revisora humana, con otros recursos y método de captura. Las últimas repeticiones en DockerDevelop usaron un único juez IA sobre las respuestas disponibles. Son procedimientos distintos, separados en el apartado~\ref{h-0}; sus puntuaciones no se suman ni demuestran superioridad general de un modelo. La ventaja de coste observada también procede de ventanas diferentes y no se atribuye sólo al modelo. \cite{reevaluacion-modelo,jup133-coste,juez-infra}

\subsection{Latencia y margen de mejora}\label{g-4}

El histórico de la tabla~\ref{tab-g1} supera el objetivo de latencia. Después se amplió el límite final y se midieron una demo y seis repeticiones en DockerDevelop, recogidas por separado en el apartado~\ref{h-0}. Falta una comparación controlada que aísle el efecto del límite, del modelo o de la infraestructura. \cite{campanas-infra,demo-final}

Las vías de ajuste, todavía sin evaluar, se recogen en el apartado~\ref{i-0}.
